\documentclass[10pt]{article}
\usepackage{icmlko}

\icmlmeta{IP 파운데이션 월드모델}{970}{잰 판본이 저장소에 없었다}{앞 노트의 수를 낸 러너가 git 어디에도 없었다 --- 다시 써서 다시 돌리니 잎 3{,}644 중 3{,}640 이 비트로 살아났고, 되살린 통제로 판을 돌리니 판이 내려간다}

\begin{document}

\twocolumn[
\icmlheader

\begin{abstract}
앞 노트(969)는 산출물 \textbf{여섯 전부}에 자기 러너의 \texttt{sha256} 을 도장으로 박았다.
\textbf{그 판본은 git 어디에도 없다} --- 객체 저장소의 blob \textbf{79{,}429 개 전량}을 훑어
일치가 \textbf{0} 이다. 969 의 사전등록 F3 는 \emph{「산출물 sha $\neq$ 커밋 blob sha 면 실패」}인데
969 는 \textbf{F3 을 \textbf{OK} 로 적었다}. 969 는 같은 사이클에서 \emph{남의} 산출물 여섯을
그 자로 대조했고 \textbf{자기 러너만 안 봤다}. 이 노트는 그것을 고친다 --- 그러나
\textbf{복원이 아니다}. 없어진 바이트는 되찾을 수 없으므로 \textbf{다시 쓰고 다시 돌렸다}.
없어진 것은 둘이었다: \textbf{§H 사료}(969 의 헤드라인인 고고학을 낸 코드)와
\textbf{§E 전후 대조}(P7 의 자리 --- 앞선 비평은 §H 만 짚었다). 커밋본으로 여섯을 다시 갈고
\textbf{잎 3{,}644 를 전량 대조}하니 \textbf{3{,}640 이 비트 동일}이고 다른 넷은 전부
\textbf{러너 자기 줄 번호}였다. \textbf{수는 살아 있었다 --- 그러나 그것은 사후에 아는 것이고,
대조하기 전까지 그 노트의 어느 수도 인용할 수 없었다.}
곁으로 다섯을 정정한다. \textbf{(i)} 「Holm 통과 불변은 재 봐야 아는 것이었다」는 과장이다 ---
두 설정 상수가 닿는 도메인은 \textbf{10 중 5} 이고, 통과 집합 셋 중 둘은
\textbf{$\Delta\rho = 0.000000$ 으로 구조상 동일}이며 나머지 하나도 문턱에서 \textbf{70--277 SE} 다.
\textbf{그 검사에는 검정력이 없었다.} \textbf{(ii)} 셋째 통과는 문턱에서 \textbf{1.88 SE} 이고
969 자신의 기계가 \texttt{씨앗을 바꾸면 뒤집힐 수 있나: true} 로 표시했다.
\textbf{(iii)} 개별 $\Delta\rho$ 를 \textbf{두 이름 중 큰 쪽}으로 골라 적었다 --- 주 이름으로는
애니가 $0.0072$ 가 아니라 $-0.000598$ 이다. \textbf{(iv)} 「0.5 채움은 층이 셋」은 세 파일만
훑은 산물이다 --- \texttt{lab/} 전량이면 \textbf{34 자리}이고, \textbf{돌려서 세면
챔피언 판에서 도는 것은 14 줄}이다. \textbf{(v)} \texttt{harness.py:241} 이
\emph{「\texttt{WIKI\_DROP} 이 떨어지는 자리」}라는 인용은 \textbf{헤드라인 도메인에서 틀리다}.
마지막으로 두 사이클이 미룬 실험을 했다: \textbf{설정 상수를 비운 채 챔피언 판 24 주행}.
$\Delta\rho = -0.002962 \pm 0.000973$ (짝지은 12 씨앗 · 3 $/$ 12).
\end{abstract}
]

\section{무엇이 없어졌나}

\claimx{969 의 수를 낸 러너는 git 어디에도 없었다 --- 사전등록 F3 가 발화했어야 했고 969 는 그것을 통과로 적었다}

969 의 산출물 여섯은 모두 \texttt{code\_stamp} 를 싣고, 그 안의
\texttt{runners/dropaudit969.py} 항목이 \texttt{cc4244b1\ldots} 다. 가지의 모든 커밋에
있는 그 파일은 \texttt{9198f932\ldots} 이고, 객체 저장소 blob \textbf{79{,}429 개} 어디에도
\texttt{cc4244b1} 은 없다. \texttt{code\_stamp} \textbf{74 파일 중 73 이 커밋과 일치}하고
\textbf{어긋난 하나가 969 자기 러너}다.

이 배치는 우연이 아니다. 969 는 같은 사이클에서 \emph{앞 노트의} 자를 수리하며
\textbf{「자는 잰 소스의 sha 를 산출물에 박는다」}는 조항을 만들었고, 그 자로 남의 산출물
여섯을 대조해 통과시켰다. \textbf{자기 러너는 그 분모에 없었다.} 도장을 찍는 자가
\emph{자기를} 분모에서 빼면 그 조건은 항진명제가 된다.

\claimx{없어진 것은 §H 사료와 §E 전후 대조 둘이고, 다시 써서 다시 돌리니 잎 3,644 중 3,640 이 비트 동일이다}

커밋본을 그대로 돌려 969 산출물과 잎별로 맞춰 보면 두 구멍이 나온다.

\begin{center}\small
\begin{tabular}{lrr}
산출물 & 잎 & \textbf{다른 잎} \\
\hline
\texttt{out969\_wiring} & 227 & \textbf{0} \\
\texttt{out969\_drop} & 2{,}440 & \textbf{0} \\
\texttt{out969\_prop} & 789 & \textbf{0}$^\dagger$ \\
\texttt{out969\_game} & 67 & \textbf{0} \\
\texttt{out969\_metacmp} & 121 & \textbf{4}$^\ddagger$ \\
\hline
합계 & \textbf{3{,}644} & \textbf{4} \\
\end{tabular}
\end{center}

\noindent{\footnotesize $^\dagger$ §E 를 다시 쓴 뒤. 커밋본은 \texttt{prop} 단계에서 §C·§D 만 내고
\textbf{§E 를 아예 안 냈다}. \quad $^\ddagger$ 넷 다 \texttt{dropaudit969.py} 의
\textbf{자기 줄 번호}다(\texttt{:893}$\to$\texttt{:910} 등) --- 없어진 코드를 써 넣었으니
줄이 밀리는 것은 불가피하고 수가 아니다.}

\medskip
§H 는 이 노트의 헤드라인인 \emph{고고학}을 낸 코드다 --- \texttt{archive/pre-github} 에서
설정 상수의 도입 커밋을 찾고, 원장의 「\ldots 잘라내기」 항목과 노트 350 의 사전등록을 읽어
\emph{「이 상수는 왜 있나」} 에 답한다. §E 는 P7 의 자리다. \textbf{둘 다 없는 채로 커밋되어
있었고, 그 사실은 산출물을 다시 돌리기 전에는 보이지 않는다.}

\section{곁으로 정정하는 다섯}

\claimx{「Holm 통과 불변」은 재 봐야 아는 것이 아니었다 --- 그 강건성 검사에는 검정력이 거의 없었다}

969 는 설정 상수를 비우고 다시 재도 Holm 통과 수가 두 이름 모두 불변임을 보이고
\emph{「이것은 재 봐야 아는 것이었다」}고 적었다. §E 를 다시 읽으면 그렇지 않다.
두 상수가 닿는 도메인은 \textbf{10 중 5}(도서 · 시장팝업 · 애니 · 웹툰 · 팝업)이고
\textbf{나머지 다섯은 $\Delta\rho = 0.000000$ 으로 구조상 동일}이다. 그런데 통과 집합은
\textbf{세계애니 · 애니 · 게임}이고 그중 \textbf{세계애니와 게임은 안 닿는 쪽}이다.
움직일 수 있는 유일한 통과인 애니도 문턱에서 \textbf{70.53 SE}(들뜸) ·
\textbf{276.8 SE}(긴 띠) 떨어져 있다. \textbf{통과 수는 달리 나올 수가 없었다.}

전후를 싣는 것은 여전히 옳다. \textbf{다만 「불변이었다」를 증거로 쓰려면 그 검사가
무엇을 뒤집을 수 있었는지를 같이 적어야 한다.}

\claimx{셋째 통과는 문턱에서 1.88 SE 이고, 969 자신의 기계가 「씨앗을 바꾸면 뒤집힌다」로 표시했다}

969 의 초록은 \emph{「문턱 근처가 아니어서 결론이 안 바뀐 것」}이라 적었다. 사전등록이 정한
50{,}000 뽑기에서 게임/들뜸은 $p = 0.0056199$ · Holm 문턱 $0.00625$ · MC SE $0.0003343$ 이라
\textbf{1.88 SE} 이고, 같은 산출물의 필드가 \texttt{씨앗을 바꾸면 뒤집힐 수 있나: true} 다.
\textbf{그 통과가 「3」을 만드는 통과다.} 969 는 별도로 200{,}000 뽑기 게임 팔을 돌려 그 결정을
굳혔지만(4.44 SE 통과), \textbf{그것은 다른 분모의 다른 주행이고 명제 절의 「통과 3」은
50{,}000 뽑기의 수다.}

\claimx{$\Delta\rho$ 를 두 이름 중 큰 쪽으로 골라 적고 어느 이름인지 밝히지 않았다}

969 는 개별 $\rho$ 의 이동을 \emph{「팝업 $0.0676$ · 도서 $0.0496$ · 시장팝업 $0.0138$ ·
애니 $0.0072$ · 웹툰 $0.0003$」} 로 적었다. §E 는 \textbf{두 이름을 갈라} 낸다:

\begin{center}\small
\begin{tabular}{lrr}
도메인 & \textbf{긴 띠 $\Delta\rho$(주 이름)} & 들뜸 $\Delta\rho$ \\
\hline
애니 & $-0.000598$ & $+0.007213$ \\
\end{tabular}
\end{center}

\noindent 969 가 인용한 수는 \textbf{둘 중 절댓값이 큰 쪽}이다. 사전등록이 정한 주 이름은
\textbf{긴 띠}이고, 그 이름으로는 애니가 $-0.000598$ 이라 969 의 예측(「애니가 $0.01$ 이상
움직인다」)이 \textbf{한 자릿수 더 크게 틀렸다}.

\claimx{「0.5 채움은 층이 셋」은 세 파일만 훑은 산물이다 --- 돌려서 세면 챔피언 판에서 도는 것은 14 줄이고, harness.py:241 인용은 헤드라인 도메인에서 틀리다}

969 의 그 표는 자기 러너가 \texttt{wikiaxes.py}·\texttt{harness.py}·\texttt{forms.py}
\textbf{세 파일만 하드코딩으로 grep} 한 결과다 --- \texttt{lab/loop.py} 는 \texttt{0.5} 로
한 번도 안 훑었다. \texttt{lab/} 전량으로 넓히면 채움 호출은 \textbf{34 자리 ·
19 파일}, \texttt{0.5} 리터럴 전량은 \textbf{94 자리 · 34 파일}이다.

정적 세기는 \emph{「있다」}만 말한다. \texttt{np.full}/\texttt{np.where} 를 감싸
채움값이 $0.5$ 인 호출의 \textbf{줄}을 세는 파수병을 달고 자료 짓기와 챔피언 한 주행을
덮으면 \textbf{14 / 34 만 실제로 돈다}. 상위 몇 줄:
\texttt{forms.py:163}(864회) · \texttt{harness.py:241}(289) ·
\texttt{trendaxes.py:54}(80) · \texttt{calaxes.py:127}(66) · \texttt{wikiaxes.py:192}(40) ·
\textbf{\texttt{loop.py:612}(31)} · \textbf{\texttt{loop.py:683}(19)}.

\textbf{그 마지막 둘이 핵심이다.} 헤드라인 도메인(팝업)은 \texttt{harness.py:241} 을
\textbf{안 지난다} --- 그 도메인이 \texttt{harness.PRIMARY} 라서 뒤에 오는 함수가 블록을
\textbf{통째로 다시 짓고}, 살아남는 $0.5$ 는 \texttt{loop.py:612} 에서 온다. 969 가 감사한
18 칸 중 \textbf{6 칸}과 검색 상수 효과의 \textbf{100\%} 가 그 줄을 비껴간다.
\textbf{문서대로 \texttt{:241} 에만 계측기를 달면 주 도메인을 조용히 놓친다.}

곁: \texttt{forms.py:166} 은 챔피언 판에서 \textbf{0 회} 돈다 --- 축 목록이 전 도메인
이름의 교집합이라 그 \texttt{else} 가지가 \textbf{원리상 도달 불가}다.

\section{두 사이클이 미룬 실험}

\claimx{설정 상수를 비우면 축 이름이 안 바뀐다 --- 그래서 「되살린 통제로 판을 돌려라」는 물을 수 있는 물음이었고, 답은 「되살리면 판이 내려간다」 다}

앞 노트는 이 실험의 전제 조건으로 \emph{「\texttt{AXIS\_MODE="common"} 의 축 붕괴를 먼저
풀어야 한다」}를 걸어 두었다. \textbf{그 전제는 틀렸다.} 축 붕괴는 \emph{열을 뺄 때} 생긴다 ---
도메인마다 \emph{다른} 열을 빼면 교집합이 무너진다. \textbf{설정 상수를 비우는 것은 열을
하나도 안 뺀다}: 그 상수는 축 dict 에서 도메인을 빼고, 그러면 판 조립기가 \textbf{이름을
남긴 채 값만 $0.5$ · 마스크 $0$ 으로} 채운다. 실측으로 현행과 비움에서 \textbf{도메인별 축
이름 수가 완전히 같다}(11 도메인 36 · 게임 40 · 교집합 36).

그래서 그대로 돌렸다 --- 짝지은 12 씨앗 $\times$ 2 팔 $=$ \textbf{24 주행}:

\begin{center}\small
\begin{tabular}{lr}
기준선(현행) & $0.470343$ \\
처리(두 상수 비움) & $0.467380$ \\
\hline
\textbf{$\Delta\rho$(짝지은 C$-$B)} & $\mathbf{-0.002962}$ \\
짝 SD & $0.003370$ \\
\textbf{짝 SE} & $\mathbf{0.000973}$ \\
$|\Delta\rho| \div$ 짝 SE & $3.04$ \\
이긴 씨앗 & 3 $/$ 12 \\
\end{tabular}
\end{center}

\noindent 기준선은 정본 $0.47034$ 를 재현한다. \textbf{채택 크기는 이 저장소가 「못 정했다」로
남긴 양이므로 판정은 $\Delta\rho \pm$ SE 로만 적는다.}

\textbf{되살리면 판이 내려간다.} 이것은 앞 노트가 캔 사료와 부호가 맞는다 --- 노트 350 은 \emph{빼는} 방향으로 판이 올랐다고 적었고 그래서 뺐다. \textbf{그러므로 「그 상수를 진짜로 지울 것인가」의 답은 「지우면 예측이 나빠진다」다.} \emph{주의} --- 분모가 다르다 --- 노트 350 은 위키 \emph{넷}만 뺐고, 이 팔은 위키 다섯과 검색 하나를 한꺼번에 되살린다. \textbf{그러나 방향은 갈리지 않는다.} \textbf{그리고 이것이 명제를 살린다}: 통제로 쓰기에 옳은 열이 판에서는 해로운 열이라는 것은 모순이 아니다 --- \textbf{두 물음이 다르다}.

\section{못 한 것}

\textbf{(1)} 없어진 §H·§E 를 \emph{바이트로} 복원하지 못했다 --- 원본이 git 어디에도 없다.
잎값이 같다는 것은 \emph{다시 쓴 코드가 같은 수를 낸다}는 뜻이고 \emph{같은 코드였다}는 뜻이
아니다. \textbf{(2)} 앞 노트가 자기 신고한 항진명제(배선 W6)를 \textbf{또 안 고쳤다} ---
찾았고 게시했고 두 사이클째 남아 있다. \textbf{(3)} 들뜸을 유보로 갈라 \emph{예측 성능}으로
재지 못했다 --- 이 노트의 모든 $\rho$ 는 여전히 \textbf{연관}이다. \textbf{(4)} 파수병은
\texttt{np.full}/\texttt{np.where} 만 본다 --- \texttt{v[\textasciitilde ok] = 0.5} 꼴의
직접 대입 채움은 못 본다. \textbf{「안 돈 줄」은 「원리상 안 돈다」가 아니라 「이 주행에서 안
돌았다」다.} \textbf{(5)} \textbf{판 팔의 도메인별 $\Delta$ 를 안 남겼다} --- 팔을 짠
함수가 씨앗마다 묶음 점수만 저장하고 도메인별 점수를 버렸다. 앞 노트의 판 팔은 그것을
실었는데 이 노트는 못 싣는다. 다시 돌리면 30 분이라 이 사이클에서는 안 돌렸다 ---
\textbf{없는 것을 있는 것처럼 적지 않는다.}

\end{document}
